\documentclass{ou}
\begin{document}

\thispagestyle{empty}
\begin{center}
\shortstack{\FGAVO}\vspace{4\baselineskip}
\textbf{РЕФЕРАТ}\\
\setstretch{2}
Информационные системы и~технологии\\
\setstretch{1}
Автоматизированные информационные системы и~системы управления\vspace{2\baselineskip}\end{center}
\begin{flushleft}
\doublespacing
\begin{tabularx}{\textwidth}{@{} p{0.4\textwidth} l l @{}}
Руководитель & ст. преподаватель & И. Р. Нигматуллин\\
Студенты     &                   &                  \\
\multirow{-2.27}{\hsize}{\singlespacing Назначение и основные характеристики автоматизированных информационных систем}
             & ГФ25-01Б, 152537684 & О. И. Бочанова\\
             & ГФ25-01Б, 152539986 & М. К. Гинтер\\
             & ГФ25-01Б, 152541471 & А. О. Безруких\\
             & ГФ25-01Б, 152536796 & А. В. Корюкова\\
\multirow{-2.27}{\hsize}{\singlespacing Архитектура и~ключевые компоненты автоматизированных систем управления}
             & ГФ25-01Б, 152540046 & А. Д. Балцевич\\
             & ГФ25-01Б, 152537659 & В. Е. Зайцева\\
             & ГФ25-01Б, 152536592 & А. Ю. Бирюкова\\    
\multirow{-2.27}{\hsize}{\singlespacing Жизненный цикл и~этапы разработки автоматизированных систем}
             & ГФ25-01Б, 152536115 & С. В. Идрисова\\
             & ГФ25-01Б, 152540131 & Е. Кичко\\
             & ГФ25-01Б, 152539526 & К. Н. Аркадьева\\
             & ГФ25-01Б, 152536056 & В. А. Карасова\\
\end{tabularx}
\end{flushleft}
\vspace*{\fill}
\begin{center}
\city
\end{center}
\newpage
\thispagestyle{empty}
\setstretch{1}
Продолжение титульного~листа реферата по~теме «Автоматизированные информационные системы и~системы управления»
\begin{flushleft}
\doublespacing
\begin{tabularx}{\textwidth}{@{} p{0.4\textwidth} l l @{}}
Студенты     &                     &         \\
\multirow{-2.27}{\hsize}{\singlespacing Современные технологии и~программные решения для~автоматизации управления}
             & ГФ25-01Б, 152540765 & А. В. Гайворонская\\
             & ГФ25-01Б, 152540753 & М. И. Бобылев\\
             & ГФ25-01Б, 152539850 & Е. А. Аксаментова\\
\multirow{-2.27}{\hsize}{\singlespacing Вопросы безопасности и~устойчивости функционирования автоматизированных систем}
             & ГФ25-01Б, 152541294 & Э. Р. Махмудов\\
             & ГФ25-01Б, 152539829 & В. Д. Бобкова\\
\multirow{-1.27}{\hsize}{\singlespacing Примеры применения автоматизированных систем в~разных сферах деятельности}
             & ГФ25-01Б, 152537671 & А. Е. Корикова\\
             & ГФ25-01Б, 152536528 & Д. А. Васильева\\
             & ГФ25-01Б, 152540016 & И. Н. Морозова\\        
\end{tabularx}
\end{flushleft}
\newpage
{
\sloppy
\tableofcontents
}
\newpage
\onehalfspacing
\section{Назначение и основные характеристики автоматизированных информационных систем}
В современном мире информационные технологии играют ключевую роль в развитии бизнеса, государственного управления и различных сферах жизнедеятельности общества. Особенно важным инструментом повышения эффективности и конкурентоспособности являются автоматизированные информационные системы (АИС). Они представляют собой сложные комплекса человеко-машинных средств, предназначенные для сбора, хранения, обработки, анализа и распространения информации для оптимизации управленческих и производственных процессов.
\subsection{Понятие и структура автоматизированных информационных систем}
Автоматизированная информационная система -- это совокупность программного обеспечения, аппаратных средств и методик, объединенных в единую архитектуру с целью автоматизации функций сбора, хранения, обработки и распространения данных. Основными компонентами АИС являются:
\begin{enumerate}[leftmargin=0em, itemindent=3.75em, nosep, label=\arabic*)]
\item Когнитивно-обрабатывающие модули: базы знаний, математические модели, методы анализа данных;
\item Телекоммуникационные средства -- сети передачи информации, обеспечивающие связь между различными элементами системы;
\item Пользовательский интерфейс -- средства взаимодействия человека с~системой, включающие системы визуализации, обработки команд;
\item Обработка и хранение данных -- базы данных, информационные архивы, системы резервного копирования;
\item Управляющие компоненты -- модули прогнозирования, принятия решений, моделирования сценариев.
\end{enumerate}

Такая интеграция позволяет организации эффективно управлять большими объемами данных и принимать обоснованные управленческие решения в короткие сроки.
\subsection{Назначение автоматизированных информационных систем}
Основная функция АИС -- создание единого информационного пространства, которое позволяет руководителям, специалистам и операторам получать актуальную и достоверную информацию в любой момент времени.

Это достигается за счет: автоматического сбора информации с~различных источников, централизованного хранения данных, формирования отчетов, автоматизированной обработки и анализа информации, рекомендаций и~сценариев дальнейших действий, поддержки принятия решений в условиях высокой неопределенности и постоянных изменений.

АИС способствуют снижению времени реагирования на внешние и~внутренние риски, автоматизации рутинных процессов, повышения точности и качества управления.

Основные сферы применения автоматизированных информационных систем:
\begin{enumerate}[leftmargin=0em, itemindent=3.75em, nosep, label=\arabic*)]
\item Промышленность и производство;
\item Государственное управление и муниципальные услуги;
\item Бизнес и коммерция;
\item Архивистика и документооборот;
\item Безопасность и профилактика рисков;
\item Технические характеристики и архитектура АИС.
\end{enumerate}

В промышленности АИС собирают данные параметров технологических процессов, контролируют работу оборудования, автоматизируют управление производственными линиями. Системы промышленной автоматизации позволяют отслеживать температуру, расход материалов и~автоматически регулировать режим работы оборудования, что повышает качество продукции, снижает издержки и предотвращает аварийные ситуации.

В государственных структурах АИС применяются для автоматизации учета граждан, ведения архивов, финансового контроля, предоставления государственных услуг в электронном виде. Это обеспечивает прозрачность, сокращает бюрократию и повышает уровень обслуживания населения.

В сфере бизнеса АИС автоматизируют учет финансовых операций, управляют ресурсами, контролируют цепочки поставок, реализуют системы электронной коммерции, автоматизируют складское хозяйство. Это позволяет повысить скорость обработки заказов, снизить операционные расходы и~увеличить прибыль.

Современные системы позволяют оцифровывать и структурировать большие объемы документации -- архивные дела, договоры, юридические документы. Их использование делает доступным поиск необходимой информации без необходимости физического перемещения документов, ускоряет обработка и хранение информации.

Интеграция АИС в системы информационной безопасности позволяет реализовать механизмы контроля доступа, шифрования данных, аудита и~мониторинга. В критических учреждениях такие системы обеспечивают защиту данных и минимизируют последствия кибератак и технических сбоев.

Современные автоматизированные системы строятся по~трехуровневой схеме. Это помогает сделать их более эффективными, надежными и легкими для расширения.

Первый уровень -- это устройства на месте работы. Там находятся датчики, которые измеряют параметры окружающей среды или оборудования, например, температуру или давление. Также есть механизмы, которые могут автоматически включать или выключать оборудование, например, вентиляторы или клапаны, в зависимости от собранных данных. Эти устройства собирают информацию в реальном времени, то есть сразу же, когда что-то происходит.

Второй уровень -- это системы, которые собирают все эти данные и~обрабатывают их. Здесь работают компьютеры и специальные программы, что позволяют контролировать работу оборудования, управлять процессами и обеспечивать автоматизм. Также здесь есть устройства, что передают информацию между полевым уровнем и высшими уровнями управления.

Третий уровень -- это верхняя часть системы, предназначенная для~человека. Там есть программы и интерфейсы, через которые руководитель или специалист может смотреть все происходящее, получать отчёты и~принимать решения. Этот уровень помогает понять, что происходит на~производстве, и скорректировать работу системы, если потребуется.

Такая структура позволяет система работать очень надежно. В~случае сбоя она может автоматически переключиться на резервные части и~продолжать работу. Также систему можно расширять -- например, подключать новые датчики или управлять новыми процессами без необходимости менять всю систему целиком.

Важно, чтобы данные хранились и передавались безопасно; для этого используют специальные механизмы защиты. И система должна быстро обрабатывать большие объемы информации, чтобы реагировать быстро и~точно, особенно в тех ситуациях, где важна мгновенная реакция.

Чтобы повысить эффективность и сделать систему еще умнее, используют современные технологии. Например, облачные сервисы позволяют хранить данные удаленно, а искусственный интеллект помогает анализировать информацию и предсказывать будущие события.
\subsection{Экономическая эффективность и перспективы внедрения}
Внедрение автоматизированных информационных систем позволяет сократить издержки на обработку информации, повысить качество продукции и услуг, оптимизировать управленческие процессы. Расчетные сроки окупаемости варьируются от нескольких месяцев (для малого бизнеса) до~двух-трех лет для крупных предприятий.

Перспективы развития АИС связаны с внедрением искусственного интеллекта, автоматизацией аналитики, повышением уровня защищенности информационных систем и использованием новых технологий для интеграции различных элементов инфраструктуры.
\newpage
\section{Архитектура и~ключевые компоненты автоматизированных систем управления}
Автоматизированная система управления (АСУ) -- это комплекс технических средств, программного обеспечения, математических методов и~персонала, предназначенный для~сбора, обработки, анализа данных и~управления объектом.

Такие системы используются в~промышленности, энергетике, транспорте, логистике и~других сферах.

\subsection{Архитектура}
Классически АСУ строится по~трёхуровневой архитектуре:
\begin{enumerate}[leftmargin=0em, itemindent=3.75em, nosep, label=\arabic*)]
\item Нижний уровень -- уровень взаимодействия с~реальными физическими процессами. Работает в~реальном времени;
\item Средний уровень -- отвечает за~управление процессом. На~этом уровне реализуются алгоритмы стабилизации процесса, защиты оборудования и~автоматизации операций.
\item Верхний уровень -- уровень визуализации и~анализа.
\end{enumerate}

Компоненты и функции нижнего уровня:
\begin{itemize}[label=-, leftmargin=0em, itemindent=3.35em, nosep]
\item датчики (температура, давление, расход, скорость, уровень и~т.д.);
\item исполнительные механизмы (вентили, клапаны, двигатели);
\item системы сбора данных (AI, DI модули, контроллеры ввода/вывода);
\item измерение параметров процесса;
\item передача данных на~верхние уровни;
\item выполнение команд управления.
\end{itemize}

Компоненты и функции среднего уровня: 
\begin{itemize}[label=-, leftmargin=0em, itemindent=3.35em, nosep]
\item ПЛК (программируемые логические контроллеры);
\item микроконтроллеры;
\item PAC-контроллеры;
\item устройства связи (шины Fieldbus, Modbus, Profibus);
\item сбор данных с~датчиков;
\item выполнение логики управления;
\item передача данных в~SCADA;
\item локальное принятие решений (включить/отключить устройство и~т.п.).
\end{itemize}

Компоненты и функции высшего уровня: 
\begin{itemize}[label=-, leftmargin=0em, itemindent=3.35em, nosep]
\item SCADA-системы (Supervisory Control And Data Acquisition);
\item серверы баз данных;
\item HMI (Human Machine Interface) панели;
\item рабочие места операторов и~инженеров;
\item визуализация данных в~реальном времени (мнемосхемы);
\item архивирование параметров (trend logs);
\item управление технологией через интерфейс оператора;
\item анализ состояния оборудования;
\item удалённый доступ;
\item диагностика работы системы.
\end{itemize}
\subsection{Взаимодействие компонентов}
АСУ работает как единый организм:
\begin{enumerate}[leftmargin=0em, itemindent=3.75em, nosep, label=\arabic*)]
\item Датчики измеряют параметры;
\item Контроллеры получают данные и~выполняют алгоритм управления;
\item Исполнительные механизмы выполняют команды контроллера;
\item SCADA визуализирует, анализирует и~логирует данные;
\item Оператор контролирует процесс и~вмешивается при~необходимости;
\item Система связи обеспечивает обмен данными;
\item Базы данных сохраняют историю для~анализа.
\end{enumerate}
\subsection{Требования к~архитектуре}
Архитектуру создает команда специалистов. Основным требованием к~главному (системному) архитектору является знание предметной области и~знание технических характеристик аппаратных и~программных средств, используемых для~построения системы. При~построении архитектуры должны быть заложены следующие свойства будущей автоматизированной системы:
\begin{itemize}[label=-, leftmargin=0em, itemindent=3.35em, nosep]
\item тестируемость (возможность установления факта правильного функционирования);
\item диагностируемость (возможность нахождения неисправной части системы);
\item ремонтопригодность (возможность восстановления работоспособности за~минимальное время при~экономически оправданной стоимости ремонта);
\item надежность (например, путем резервирования);
\item простота обслуживания и~эксплуатации (минимальные требования к~квалификации и~дополнительному обучению эксплуатирующего персонала);
\item экономичность (эффективность в~процессе функционирования);
\item наращиваемость (возможность увеличения автоматизированной системы при~увеличении размера объекта автоматизации).
\end{itemize}
\newpage
\section{Жизненный цикл и~этапы разработки автоматизированных систем}
Автоматизированная система (АС) -- это комплекс программных, технических и~информационных средств, предназначенных для~автоматизации деятельности в~определенной предметной области. Разработка и~внедрение АС требует систематического подхода, который реализуется через жизненный цикл системы. Жизненный цикл АС – это период времени, который начинается с~момента возникновения потребности в~системе и~заканчивается полным ее выводом из~эксплуатации.
\subsection{Понятие и~модели жизненного цикла АС}
Жизненный цикл АС включает последовательность взаимосвязанных процессов, этапов и~задач, направленных на~создание и~сопровождение системы. Существуют различные модели жизненного цикла, которые описывают эту последовательность, но~все они преследуют общую цель -- обеспечить управляемый и~эффективный процесс разработки и~эксплуатации АС.

Наиболее распространенные модели жизненного цикла:
\begin{itemize}[label=-, leftmargin=0em, itemindent=3.35em, nosep]
\item каскадная модель (водопадная);
\item инкрементная модель;
\item спиральная модель;
\item Agile-модели (Scrum, Kanban).
\end{itemize}

Каскадная модель предполагает последовательное выполнение этапов разработки, где каждый этап завершается перед переходом к~следующему. Требования к~системе должны быть четко определены заранее. Преимущества: простота управления, четкая структура. Недостатки: сложность внесения изменений на~поздних этапах, не~подходит для~проектов с~неопределенными требованиями.

В инкрементной модели система разрабатывается и~внедряется по~частям (инкрементам). На~каждом этапе создается работоспособная, хотя и~неполная, версия системы. Преимущества: возможность ранней демонстрации функциональности, гибкость к~изменениям. Недостатки: требует четкой координации между~инкрементами.

Спиральная модель сочетает итеративную разработку с~элементами управления рисками. Каждый виток спирали представляет фазу разработки, включающую планирование, анализ риска, разработку, оценку результатов. Преимущества: управление рисками на~каждом этапе, возможность внесения изменений. Недостатки: требует высокой квалификации разработчиков и~управления рисками.

Agile-модели основаны на~итеративной разработке с~акцентом на~гибкость, взаимодействие с~заказчиком и~быструю адаптацию к~изменениям. Преимущества: быстрая обратная связь, гибкость, ориентированность на~результат. Недостатки: требует высокой вовлеченности заказчика, сложнее в~управлении крупными проектами.

Выбор модели жизненного цикла зависит от~специфики проекта, требований к~системе, доступных ресурсов и~опыта команды разработчиков.
\subsection{Этапы разработки автоматизированных систем}
Несмотря на~различия в~моделях жизненного цикла, можно выделить общий набор этапов, которые присущи большинству проектов по~разработке АС:
\begin{itemize}[label=-, leftmargin=0em, itemindent=3.35em, nosep]
\item анализ требований;
\item проектирование системы;
\item реализация системы (программирование);
\item тестирование системы;
\item внедрение системы;
\item эксплуатация и~сопровождение системы.
\end{itemize}

При анализе требований нужно определить потребности заказчика, собрать и~провести анализ требований к~системе, сформулировать цели и~задачи АС. На этом этапе проводится интервью с~заказчиком и~пользователями, анализ существующих систем и~бизнес-процессов, разработка спецификации требований (SRS - Software Requirements Specification). Результатом является спецификация требований, технико-экономическое обоснование (ТЭО) проекта.

Цель проектирования системы -- разработка архитектуры системы, определение компонентов, интерфейсов, структуры данных и~алгоритмов. Разрабатывается концептуальная, логическая и~физическая модели системы, проектируются базы данных, разрабатывается пользовательский интерфейс (UI/UX). Результат -- проектная документация, макеты интерфейсов, схема базы данных.

Реализация системы должна предоставить разработку программного кода на~основе проектной документации. Действия -- написание кода, отладка, модульное тестирование компонентов системы. Итог -- исходный код программного обеспечения.

Тестирование системы нужно для выявления ошибок и~несоответствий в~программном коде и~работе системы. Разрабатываются тестовые сценарии, проводятся тестирования различных видов (модульное, интеграционное, системное, приемочное). Результат -- отчеты об~ошибках, исправленный программный код.

Целью внедрения системы является развертывание системы в~рабочей среде, обучение пользователей, миграция данных и~интеграция с~другими системами. Осуществляется установка программного обеспечения, настройка системы, обучение персонала, перенос данных из~старых систем, запуск системы в~эксплуатацию. Результат -- готовая к~работе АС.

Эксплуатация и~сопровождение системы обеспечивает поддержку работоспособности системы, устранение возникающих ошибок, внесение изменений и~улучшений в~соответствии с~изменяющимися потребностями пользователей. При этом этапе происходит мониторинг работы системы, исправление ошибок, обновление программного обеспечения, адаптация системы к~новой среде, разработка новых функций. Итог -- стабильно функционирующая и~развивающаяся АС.
\subsection{Инструменты и~методы разработки АС}
Для эффективной разработки АС используются различные инструменты и~методы, которые позволяют автоматизировать процессы, повысить качество и~снизить затраты:
\begin{enumerate}[leftmargin=0em, itemindent=3.75em, nosep, label=\arabic*)]
\item CASE-средства -- инструменты, предназначенные для~автоматизации этапов разработки АС (например, AllFusion ERwin Data Modeler, Rational Rose);
\item Системы управления проектами (PMS) -- инструменты для контроля, организации и~планирования выполнения задач в~проекте (например, Microsoft Project, Jira);
\item Системы контроля версий (VCS) -- инструменты для~управления изменениями в~исходном коде (например, Git, SVN);
\item Инструменты для~автоматизированного тестирования -- инструменты для~автоматизации процессов тестирования, что позволяет сократить время и~повысить качество тестирования. (например, Selenium, JUnit);
\item Методологии разработки (Agile, Scrum, Kanban) -- наборы принципов, практик и~инструментов для~организации процесса разработки.
\end{enumerate}

Жизненный цикл АС представляет сложный и~многоэтапный процесс, который требует систематического подхода и~применения современных методов и~инструментов разработки. Выбор подходящей модели жизненного цикла, тщательное планирование, эффективное управление рисками и~тесное взаимодействие с~заказчиком являются ключевыми факторами успешной разработки и~внедрения автоматизированной системы.
\newpage
\section{Современные технологии и~программные решения для~автоматизации управления}
\subsection{Понятие автоматизации управления и~её роль}
Автоматизация управления -- это использование оборудования, технологий и~программного обеспечения для~выполнения управленческих функций без~непосредственного участия человека или~с~его минимальным участием. Она охватывает процессы анализа информации, планирования, контроля, распределения ресурсов и~принятия решений.

Основные задачи автоматизации управления:
\begin{enumerate}[leftmargin=0em, itemindent=3.75em, nosep, label=\arabic*)]
\item Сокращение затрат и~повышение производительности. Автоматизация позволяет снизить количество ручной работы и~ускорить выполнение операций;
\item Повышение качества управления;
\item Минимизация человеческого фактора. Исключение ошибок и~повышение надёжности;
\item Централизация данных и~процессов. Единые информационные платформы позволяют принимать решения на~основе актуальной информации.
\end{enumerate}

Пример применения автоматизации управления: крупная транспортная компания внедрила автоматизированную систему распределения грузов между~водителями. Алгоритм автоматически рассчитывает оптимальные маршруты, учитывая трафик, загрузку автомобилей и~срочность заказов. В~результате удалось сократить время доставки на~15\% и~уменьшить расход топлива на~10\%.
\subsection{Современные технологии автоматизации управления}
В последние годы появилось множество новых технологий, радикально изменивших подход к~автоматизации.

Облачные платформы позволяют хранить данные и~запускать приложения на~удалённых серверах. Это снижает стоимость инфраструктуры и~повышает масштабируемость. Среди преимуществ: отсутствие необходимости содержать собственные серверы, быстрый доступ к~данным из~любой точки, высокая гибкость.

Пример использования: компания среднего бизнеса перенесла свою систему бухгалтерского учёта в~облако Microsoft Azure. Благодаря этому такие задачи, как формирование отчётности и~подготовка документов, стали выполнять удалённые сотрудники, что повысило оперативность работы.

Big Data позволяют анализировать огромные массивы информации в~режиме реального времени, что особенно важно при~управлении сложными системами. Это открывает возможности: прогнозирования спроса, выявления закономерностей поведения клиентов, оптимизации производственных процессов.

Ритейлер использует аналитику больших данных для~прогнозирования продаж в~различных регионах. Система анализирует историю покупок, сезонность, погоду, акции конкурентов. Это позволило сократить излишние запасы на~складах на~25\%.

ИИ и~машинное обучение позволяют автоматизировать интеллектуальные управленческие функции: прогнозирование, классификацию, принятие решений, распознавание образов. К типичным задачам ИИ и машинного обучения относят: прогнозирование риска, автоматическую обработку обращений клиентов, оптимизацию маршрутов и~логистических цепочек.

Банк внедрил интеллектуальную систему скоринга, анализирующую кредитоспособность клиентов на~основе сотен параметров. Благодаря ИИ уровень просроченной задолженности снизился на~12\%.

IoT -- это сеть устройств, обменяющихся данными между~собой. В~управлении эта технология используется для~мониторинга оборудования, контроля ресурсов и~автоматизации бизнес-процессов в~режиме реального времени.

Пример: в производственном цехе датчики контролируют температуру и~вибрацию оборудования. Когда параметры выходят за~нормы, система автоматически вызывает техника и~останавливает оборудование. Это снизило количество аварийных ситуаций на~40\%.

RPA -- это программные роботы, которые выполняют рутинные операции: ввод данных, перенос информации из~одной системы в~другую, обработка заявок. Среди преимуществ: высокая скорость работы, исключение ошибок, круглосуточная доступность.

Страховая компания внедрила RPA-робота, который автоматически обрабатывает заявления о~возмещении ущерба. Если раньше сотрудник тратил 20 минут на~каждую заявку, то теперь робот выполняет эту задачу за~2 минуты без~участия человека.
\subsection{Программные решения для~автоматизации управления}
Существуют различные классы программных систем, используемых для~автоматизации.

ERP -- комплексные системы для~управления всеми ресурсами компании: финансами, персоналом, логистикой, производством и~складом.

Популярные ERP:
\begin{itemize}[label=-, leftmargin=0em, itemindent=3.35em, nosep]
\item SAP ERP;
\item «1С:ERP»;
\item Microsoft Dynamics 365;
\item Oracle ERP Cloud.
\end{itemize}

Промышленное предприятие внедрило SAP ERP для автоматизации закупок, продаж и~планирования производства. В~результате затраты на~складские операции снизились на~18\%, а~время на~формирование отчетности -- в~4 раза.

CRM используются для~управления взаимоотношениями с~клиентами: продажи, маркетинг, сервис.

Функции:
\begin{itemize}[label=-, leftmargin=0em, itemindent=3.35em, nosep]
\item ведение базы клиентов;
\item автоматизация продаж;
\item контроль работы менеджеров;
\end{itemize}

Популярные решения:
\begin{itemize}[label=-, leftmargin=0em, itemindent=3.35em, nosep]
\item Bitrix24;
\item AmoCRM;
\item Salesforce;
\item HubSpot;
\end{itemize}

Пример: компания малого бизнеса внедрила CRM Bitrix24. Благодаря автоматизации воронки продаж продажи за~полгода выросли на~20\%, а~время обработки заявки сократилось более чем вдвое.

MES предназначены для~оперативного управления производственными процессами.

Функции:
\begin{itemize}[label=-, leftmargin=0em, itemindent=3.35em, nosep]
\item контроль работы оборудования;
\item управление заказами;
\item учёт производственных операций.
\end{itemize}

Пример: завод электроники внедрил MES-систему, которая отслеживает все этапы сборки изделий в~режиме реального времени. Это снизило процент брака на~15\%.

BI-платформы используются для~анализирования данных и~визуализации показателей деятельности.

Примеры:
\begin{itemize}[label=-, leftmargin=0em, itemindent=3.35em, nosep]
\item Power BI;
\item Qlik Sense;
\item Tableau.
\end{itemize}

Возможности:
\begin{itemize}[label=-, leftmargin=0em, itemindent=3.35em, nosep]
\item построение интерактивных отчётов;
\item прогнозирование тенденций;
\item анализ KPI.
\end{itemize}

Пример: компания по~доставке еды использует Power BI для~мониторинга заказов и~работы курьеров. Это позволило оптимизировать распределение заказов и~повысить скорость доставки.

BPM-системы часто помогают моделировать, оптимизировать, автоматизировать бизнес-процессы.

Примеры:
\begin{itemize}[label=-, leftmargin=0em, itemindent=3.35em, nosep]
\item ELMA BPM;
\item ARIS;
\item Camunda.
\end{itemize}

Финансовая организация внедрила BPM-систему для автоматизации процесса выдачи займов. Срок рассмотрения заявки сократился с~5 суток до~1 дня.
\subsection{Комплексная автоматизация управления}
Современные организации используют сразу несколько решений -- ERP, CRM, BI, облака, ИИ. В~совокупности они формируют цифровую экосистему предприятия.

Интеграционные шины данных, API и~микросервисные архитектуры позволяют связать разные программные продукты.

Пример: в торговой сети ERP, CRM и~система видеонаблюдения интегрированы в~единую платформу. Это позволяет в~режиме реального времени отслеживать продажи и~поведение покупателей в~торговом зале.

Проблемы комплексной автоматизации:
\begin{itemize}[label=-, leftmargin=0em, itemindent=3.35em, nosep]
\item высокая стоимость внедрения;
\item необходимость обучения персонала;
\item проблемы совместимости ПО.
\item риски киберугроз.
\end{itemize}
\subsection{Перспективы развития технологий автоматизации}
Автоматизация управления продолжает активно развиваться.

Гиперавтоматизация подразумевает сочетание RPA, ИИ, машинного обучения и~аналитики. Системы будут самостоятельно анализировать процессы и~предлагать пути их оптимизации.

Цифровые двойники -- модели реальных объектов, позволяющие управлять ими дистанционно.

Пример: на заводах создаются цифровые двойники оборудования, что позволяет диагностировать неполадки ещё до~их возникновения.

Будущее за~системами, которые принимают решения без~участия человека (например, автономные склады Amazon).
\newpage
\section{Вопросы безопасности и~устойчивости функционирования автоматизированных систем}
Актуальность темы обусловлена внедрением автоматизированных систем (АС) во~все сферы деятельности. Под безопасностью АС понимается свойство, обеспечивающее защиту информации, обрабатываемой системой, а~также~устойчивость ее функционирования к~непреднамеренным или преднамеренным воздействиям дестабилизирующего характера. Устойчивость функционирования АС -- это способность выполнять возложенные функции в~условиях внешних дестабилизирующих воздействий и~внутренних отказов ГОСТ Р 53647.6-2012. Менеджмент непрерывности бизнеса. Требования к~системе менеджмента непрерывности бизнеса.
\subsection{Классификация угроз безопасности автоматизированных систем}
Угрозы безопасности АС подразделяются на~внутренние и~внешние. Внутренние угрозы связаны с~недостатками проектирования, реализации и~эксплуатации АС. Внешние угрозы обусловлены целенаправленными действиями злоумышленников или~воздействиями окружающей среды.

Распространенной внешной угрозой являются кибератаки, включая внедрение вредоносного программного обеспечения. Отмечается, что значительная доля компонентов АС производится в~различных странах, что может создавать дополнительные риски на~этапе поставки, ведь многие аппаратные компоненты производятся, например, в~Китае.

\subsection{Основные аспекты обеспечения безопасности АС}
Обеспечение безопасности АС является комплексной задачей и~включает следующие аспекты:

\begin{enumerate}[leftmargin=0em, itemindent=3.75em, nosep, label=\arabic*)]
\item Организационный аспект. Включает разработку политик безопасности, регламентов и~процедур управления доступом;
\item Технический (технологический) аспект. Включает применение средств защиты информации, таких как межсетевые экраны, системы обнаружения вторжений и~средства криптографической защиты;
\item Программный аспект. Заключается в~использовании лицензионного программного обеспечения, регулярном применении обновлений безопасности и~проведении аудита кода.
\end{enumerate}
\subsection{Принципы обеспечения устойчивости функционирования АС}
Устойчивость АС достигается за~счет реализации ряда принципов:
\begin{enumerate}[leftmargin=0em, itemindent=3.75em, nosep, label=\arabic*)]
\item Избыточность. Дублирование критически важных компонентов системы (серверов, каналов связи, источников питания);
\item Масштабируемость -- способность наращивать вычислительную мощность и~ресурсы без~кардинального изменения архитектуры;
\item Отказоустойчивость. Свойство системы сохранять работоспособность при~возникновении отдельных отказов ее компонентов;
\item Восстанавливаемость. Наличие механизмов быстрого восстановления работоспособности АС после сбоев, основанных на~резервном копировании данных и~репликации.
\end{enumerate}
\subsection{Взаимосвязь безопасности и~устойчивости}
Безопасность и~устойчивость функционирования АС являются взаимосвязанными понятиями. Реализация успешной кибератаки, такой как ransomware-инцидент, напрямую подрывает устойчивость функционирования АС, приводя к~остановке бизнес-процессов . В~свою очередь, недостаточная устойчивость, выражающаяся в~отсутствии резервных копий, усугубляет последствия инцидента информационной безопасности.

Таким образом, вопросы безопасности, устойчивости функционирования автоматизированных систем требуют комплексного подхода. Эффективное управление рисками в~данной области возможно только при~одновременном применении организационных, технических, программных мер, направленных как на~предотвращение инцидентов, так и~на~минимизацию их последствий. Дальнейшие исследования целесообразно направить на~разработку моделей оценки совместной безопасности и~устойчивости АС.
\newpage
\section{Примеры применения автоматизированных систем в~разных сферах деятельности}
В современных условиях автоматизация становится общедоступной в~силу широкого распространения интернет-технологий. Многочисленные онлайн-сервисы и~программные решения предоставляют возможности для~автоматизации широкого спектра повседневных и~профессиональных задач, что способствует экономии временных и~трудовых ресурсов.

Ключевым аспектом в~контексте обсуждения автоматизации является оптимизация. Внедрение автоматизированных систем позволяет повысить эффективность выполнения операций, минимизировать ошибки, связанные с~человеческим фактором, а~также увеличить скорость и~точность рабочих процессов. Задачи, на~выполнение которых ранее требовались значительные временные затраты, могут быть реализованы с~помощью автоматизации в~сжатые сроки и~с~меньшими издержками.

Фундаментальный принцип автоматизации заключается в~применении специализированных технических средств, таких как специальные механизмы и~приводы. Данные устройства, основанные на~физических и~механических принципах, преобразуют один вид энергии в~другой и~обеспечивают передачу усилий для~выполнения целевых действий. Их использование позволяет автоматизировать процессы различной степени сложности -- от~решения бытовых вопросов до~управления комплексными производственными линиями.

К основным преимуществам автоматизации относятся:
\begin{itemize}[label=-, leftmargin=0em, itemindent=3.35em, nosep]
\item Эффективность сокращение затрат времени и~труда при~одновременном повышении качества и~точности выполнения операций;
\item Производительность обеспечивает рост объемов выполняемой работы и~увеличение общих показателей производительности;
\item Точность предоставляет снижение влияния человеческого фактора, уменьшение количества ошибок;
\item Гибкость дает возможность простой адаптации и~перенастройки систем в~соответствии с~изменяющимися требованиями.
\end{itemize}
\subsection{Поддержка креативности и~стратегического развития сотрудников}
Технологии автоматизации позволяют освобождать сотрудников от~рутинных задач, предоставляя им возможность сосредоточиться на~более креативных и~стратегически важных видах деятельности.

Это способствует увеличению производительности работы и~адаптации организации к~изменяющимся условиям рынка, что является ключевым фактором конкурентоспособности.

Крупнейшие российские банки, включая Сбербанк и~ВТБ, внедрили роботизированные системы автоматизации процессов (RPA) для~обработки заявок на~кредитные продукты. Роботы осуществляют верификацию документов, проверку кредитной истории и~расчет параметров кредита, что сократило время обработки заявки с~24 часов до~15 минут. Высвобожденные сотрудники отделения кредитного анализа переведены в~департамент развития цифровых продуктов, где занимаются созданием новых финансовых сервисов и~анализом рыночных тенденций.

Сеть гипермаркетов «Лента» внедрила систему автоматической консолидации отчетности на~платформе SAP. Ежедневные операционные отчеты по~товарным остаткам, продажам и~движению товаров теперь формируются без~участия персонала. Менеджеры по~развитию, ранее затрачивавшие до~60\% рабочего времени на~составление отчетов, теперь переориентированы на~задачи анализа потребительского поведения, разработки программ лояльности, что позволило увеличить повторные покупки на~17\% за~последний квартал.
\subsection{Улучшение управления бизнес-процессами}
Современные методы автоматизации дают значительные преимущества в~управлении и~контроле компаний. Посредством автоматизированного сбора, обработки и~анализирования информации, организации получают полное представление о~своей деятельности. Это дает возможность принимать обоснованные решения, прогнозировать рыночные тенденции и~эффективно распределять ресурсы, что в~конечном итоге ведет к~увеличению прибыли.

Автоматизация не~только повышает производительность, но~и~открывает новые возможности для~анализа данных, что критически важно в~бизнесе.

ОАО «Камаз»  реализовало проект внедрения ERP-системы Oracle для~интеграции данных из~42 структурных подразделений. Система автоматически отслеживает движение деталей по~цехам, прогнозирует потребность в~материалах и~оптимизирует загрузку производственных мощностей. В~результате достигнуто:

\begin{itemize}[label=-, leftmargin=0em, itemindent=3.35em, nosep]
\item снижение уровня незавершенного производства на~23\%;
\item сокращение логистических издержек на~15\%;
\item уменьшение времени переналадки оборудования на~30\%.
\end{itemize}

Компания «Procter \& Gamble» внедрила система автоматического пополнения запасов на~основе технологий искусственного интеллекта. Система анализирует данные с~POS-терминалов 15 000 торговых точек, прогнозирует спрос с~точностью 94\% и~автоматически формирует заказы дистрибьюторам. Это позволило:
\begin{itemize}[label=-, leftmargin=0em, itemindent=3.35em, nosep]
\item оптимизировать складские запасы на~28\%;
\item сократить операционные затраты на~логистику на~22\%.
\end{itemize}

Госкорпорация «Росатом» автоматизировала процесс бюджетного управления с~помощью платформы «1С: Управление производственным предприятием». Система консолидирует данные от~120 дочерних обществ, автоматически проверяет соответствие показателей KPI и~формирует прогнозные модели. Внедрение позволило:
\begin{itemize}[label=-, leftmargin=0em, itemindent=3.35em, nosep]
\item сократить цикл бюджетного планирования с~45 до~14 дней;
\item увеличить точность прогнозирования денежных потоков на~35\%;
\item снизить количество ошибок в~финансовой отчетности на~90\%.
\end{itemize}
\subsection{Производственная сфера}
Производственная отрасль стала одной из~первых, в~которой произошла широкая интеграция автоматизации. Благодаря внедрению промышленных роботов и~систем управления, предприятия существенно увеличили эффективность производства, одновременно уменьшая затраты. Применение автоматизации ускоряет процесс создания продуктов и~улучшает их качество, снижает количество дефектов и~повышает безопасность на~рабочих местах.

Использование роботов для~выполнения операций, таких как сварка и~сборка, способствует уменьшению времени и~повышению качества производства.

Пример: концерн Volkswagen на~своем заводе в~Вольфсбурге использует полностью автоматизированный цех кузовной сборки. Линия оснащена промышленными роботами KUKA KR QUANTEC, которые выполняют более 90\% операций точечной сварки кузова. Это обеспечивает погрешность позиционирования менее 0,2 мм и~позволяет производить смену моделей на~одной линии без~остановки производства.

Системы AS/RS оптимизируют хранение и~перемещение товаров, уменьшая ручной труд и~количество ошибок

Пример: логистические центры компании Amazon используют систему роботов-перевозчиков Kiva. Роботы автоматически перемещают стеллажи с~товарами к~станциям комплектации заказов. Внедрение данной системы позволило увеличить производительность склада на~40\% и~сократить время отбора заказа с~60-75 минут до~15 минут.

Системы контроля качества оснащены сенсорами для~выявления несоответствий стандартам, что ускоряет выявление и~устранение дефектов продукции.

Пример: на~заводе компании Bosch по~производству топливной аппаратуры используются системы машинного зрения на~базе камер Cognex. Данные системы выполняют 100\% контроль микроскопических отверстий в~форсунках, выявляя дефекты размером до~5 микрон, что невозможно для~человеческого глаза. Это позволило снизить процент брака на~0,001\%.
\subsection{Роль автоматизации в~сфере услуг}
Автоматизация играет значительную роль в~улучшении сервиса и~внутренних процессов в~сфере услуг.

Внедрение технологий, как чат-боты и~системы онлайн-бронирования, позволяет организациям быстрее реагировать на~запросы клиентов, улучшая обслуживание и~снижая расходы.

Чат-боты и~виртуальные ассистенты брабатывают запросы клиентов мгновенно, сокращая ожидание и~освобождая сотрудников для~более сложных задач

Пример: крупнейший российский банк Сбербанк внедрил виртуального ассистента «Сбер». Ассистент обрабатывает до~70\% типовых запросов клиентов, таких как проверка баланса, блокировка карты или~информация о~тарифах, что позволило сократить нагрузку на~контакт-центр на~30\%.

Благодаря системе бронирования клиенты могут самостоятельно выбирать и~оплачивать услуги, что уменьшает нагрузку на~персонал и~упрощает процесс.

Пример: авиакомпания «Аэрофлот» использует автоматизированную систему управления доходами (Revenue Management System) от~PROS. Система в~реальном времени анализирует спрос, конкуренцию и~поведение клиентов, динамически изменяя стоимость авиабилетов, что обеспечивает до~5\% дополнительного годового дохода.

CRM-системы помогают лучше понимать нужды клиентов, повышая персонализацию обслуживания через сбор и~анализ клиентских данных.

Пример: компания Salesforce предоставляет платформу для~автоматизации продаж и~обслуживания клиентов. Например, она автоматически присваивает лидам рейтинг, распределяет заявки между~менеджерами по~заданным правилам и~отправляет триггерные письма, что, по~данным компании, увеличивает конверсию на~27\%.
\subsection{Инновации в~медицинской сфере: ускорение диагностики и~лечения}
Использование автоматизации в~медицине значительно изменяет подход к~диагностике и~лечению пациентов. Она позволяет быстрее ставить диагнозы, повышая точность терапии и~уменьшая рабочую нагрузку на~врачей. Современные технологии автоматизации делают процесс медицинского обслуживания более эффективным, что положительно сказывается на~качестве оказания услуг.

Роботизированные платформы для~хирургии предоставляют возможность хирургам выполнять сложные операция с~высокой точностью и~меньшей инвазивностью. Благодаря числовому управлению, эти системы выполняют действия с~микронной точностью, уменьшая риски постоперационных осложнений и~ускоряя время восстановления пациентов. Внедрение робототехники в~медицине позволяет повышать не~только точность, но~и~безопасность проводимых вмешательств.

Пример: в~Национальном медицинском исследовательском центре хирургии им. Вишневского применяется роботизированный комплекс da Vinci~Xi. Он позволяет проводить сложные кардиохирургические, онкологические операции через минимальные разрезы. Точность манипуляций системы снижает интраоперационную кровопотерю на~30\% и~сокращает срок госпитализации на~25\%.

Автоматизированные системы диагностики используют искусственный интеллект и~алгоритмы машинного обучения для~более точного анализа медицинской информации и~изображений. Эти решения быстро определяют потенциал патологий на~снимках, что позволяет врачам оперативно принимать решения и~начинать лечение.

Пример: лабораторная служба «Инвитро» использует полностью автоматизированные биохимические анализаторы Roche Cobas 8000. Одна установка способна проводить до~3000 тестов в~час с~минимальным участием лаборанта, что ускоряет получение результатов пациентом в~2 раза и~исключает ошибки при~ручном вводе данных.

Электронные медицинские записи (EMR) автоматизируют многие аспекты управления данными пациентов. Благодаря этим системам улучшается координация лечения, а~информация становится более доступной и~актуальной, что снижает вероятность ошибок. Электронные записи синхронизируются автоматически, обеспечивая надёжный обмен данными между~медицинскими учреждениями.

Пример: единая государственная информационная система в~сфере здравоохранения (ЕГИСЗ) обеспечивает ведение электронных медицинских записей. Это позволяет врачу в~любом медицинском учреждении, имеющем доступ, просмотреть полную историю болезни пациента, включая результаты анализов и~выписные эпикризы, что критически важно для~преемственности лечения.

Сегодня внедрение автоматизированных систем управления является стратегической необходимостью для~современного предприятия, стремящегося к~устойчивому развитию и~сохранению конкурентоспособности. Их универсальность и~многозадачность позволяют создавать адаптированные решения для~различных отраслей, формируя единую информационную среду для~централизованного контроля и~принятия обоснованных решений. В~конечном счете, АСУ решают задачи, недостижимые традиционными методами, обеспечивая повышение эффективности.
\newpage
{
\sloppy
\nocite{*}
\cleardoublepage
\addcontentsline{toc}{section}{Список использованных источников}
\bibliographystyle{plainurl}
\bibliography{report3}
}
\end{document}